% Folder 79, batch 07 — archivists' pages 121–140.
% Pass: Opus 5.5 (claude-opus-5-5), 2026-10-03 — first pass, unchecked against the pages by a human.
%
% His own pagination is visible on the rectos (p. 122 = 62, p. 124 = 63,
% p. 126 = 64, …); the pages between continue the text. The diagonal
% marginal additions in the left margin are dense and often only partly read.
\documentclass[11pt,a4paper]{article}
\input{../preamble/grothendieck}

\folder{79}
\batch{7}
\pages{121}{140}
\foldertitle{2-polyèdres réguliers triangulés, et modules libres de rangs 2 sur les anneaux $\Lambda=\mathbb{Z}/n\mathbb{Z}$ et $\Lambda=\mathbb{Z}$ : notes manuscrites (s.d.).}
\dating{[à partir de 1980]}
\watermark{Édition de démonstration}

\begin{document}

\page{121}\note{la page continue un argument commencé avant ce lot : $\lambda$, $P_0$, $P_1$, $A_0$, $A_1$ y sont déjà introduits.}
(45) $\forall (s,t)\in S$, on a $\{s,t\}\in A_0 \Longleftrightarrow \{s,\lambda t\}\in A_1$.

On voit que cela implique que les groupes d'automorphismes $G_0$, $G_1$ de $P_0$, $P_1$
sont égaux, et que par suite $\lambda$ \add{\ill{} $\lambda^{-1}$} \uncertain{commute aussi à} $G_1$,
et que l'ensemble $A_0$ est « transformé » de $A_1$ par l'élément $\lambda^{-1}$ du
\uncertain{normalisateur} de $G_1$ dans $\mathfrak{S}_S$ : la relation \uncertain{envisagée} est
\ill{} symétrique, elle est évidemment réflexive, et on voit tout de suite qu'elle
est également transitive : \struck{elle} c'est donc bien une \emph{relation d'équivalence} entre
structures $\pi$-polyédrales spéciales. Si \struck{\ill{}} $\pi$ est une classe de
telles structures, il lui correspond un groupe
\marginal{en diagonale, à gauche : on note l'\uncertain{ensemble} \ill{} radicalement \ill{}}

(46) $G_0$, groupe d'automorphismes des structures \uncertain{équivalentes} $P_i$
\struck{$\in$ \ill{}} ($i\in\pi$)

et un groupe

(47) $Z' = \mathrm{Centr}_{\mathfrak{S}_S}(G_0)$,

enfin, on voit que le groupe $Z'$ opère sur $\pi$, transitivement par définition

\page{122}\note{p. 62 de l'auteur.}
des autres relations d'équivalence. [C'est le moment \uncertain{précis}, d'ailleurs, de
noter la relation plausible

(48) $Z'$ est un groupe \emph{commutatif}
\marginal{en diagonale, à gauche : pour \ill{} (48) pour \uncertain{un axiome} \ill{} \uncertain{pour l'identité} \ill{}}

En effet, si $\lambda,\mu\in Z'$, et $(s,t)\in S^2$, on a :
\[
(s,t)\in A_0 \overset{\lambda\in Z'}{\Longleftrightarrow} \{\lambda s,\lambda^{-1}t\}\in A_0
\overset{\mu\in Z'}{\Longleftrightarrow} \{\mu\lambda s,\mu^{-1}\lambda^{-1}t\}\in A_0
\]
\[
\Big\Updownarrow\ \lambda\mu\in Z' \qquad
\{(\lambda\mu)s,(\lambda\mu)^{-1}t\}\in A_0
\]
\uncertain{dans une conséquence} \ill{} les extrêmes, et posant
$\mu^{-1}\lambda^{-1}t=(\lambda\mu)^{-1}t=u$ :
\[
\{u,(\lambda\mu)s\}\in A_0 \Longleftrightarrow \{u,(\mu\lambda)s\}\in A_0
\]
$\forall s,u\in S$, ce qui s'écrit
\[
A_0((\lambda\mu)s) = A_0((\mu\lambda)(s)),
\]
\struck{\ill{}} égalité des \emph{étoiles} en $(\lambda\mu)(s)$ et $(\mu\lambda)(s)$.
\uncertain{Ceci} \ill{} que \uncertain{les sommets} sont \uncertain{déterminés} par leur
étoile, \ill{} (48) — \uncertain{mais} \ill{} \uncertain{mais pas} \struck{\ill{}}
l'\uncertain{écrire} \uncertain{explicite} sans \ill{} déterminé par son étoile — (qui est
d'ailleurs \uncertain{faux} pour l'octaèdre …).]
\marginal{en diagonale, à gauche, trois lignes barrées d'un trait : c'est vrai dès que
$P_0$ n'est pas \ill{} \ill{} octaèdre \ill{}}
\marginal{en diagonale, au-dessous : cf.\ \ill{} p.~22 \ill{} \ill{}}

\page{123}
La commutativité de $Z'$ \uncertain{signalée} résulte déjà des axiomes To~3, To~4
(To~1 et To~2 sont \ill{}) ; pour un 2-polyèdre triangulé connexe, \uncertain{dont} un
sommet n'est pas déterminé par son \uncertain{étoile} \struck{(à \ill{} \ill{} ($n\geqslant 3$))}
\add{\ill{}}, \uncertain{le polyèdre} est \uncertain{pair} \ill{} $n\neq 4$ \ill{}
\uncertain{l'automorphisme} \ill{}, le groupe des \uncertain{automorphismes} est
\ill{} diédral $D_n$ ($\simeq G_0^{+}$) \ill{} le groupe $\{1,\varepsilon\}$ engendré
\ill{} $\mathfrak{S}_S$, $G_0^{+}$ est \uncertain{commutatif} \ill{}
\struck{$\{1,\varepsilon\}\times\{1,\varepsilon\}$, \ill{} 3 \ill{}}
\note{ce début de page est surchargé d'ajouts interlinéaires et biffures ; la lecture
en est très incertaine.}

réduit : $G_0^{+}\simeq D_n$ (réduit : $\{1\}$ si $n$ impair, \ill{} dans ce cas
pair) \ill{} est commutatif. Reste le cas \ill{} $n=4$, \ill{} \ill{} \ill{} l'octaèdre.
Dans le cas $G_0^{+}$ ($\simeq\mathfrak{S}_4$ \uncertain{cf.\ p.} \ill{}) qui \ill{}
\ill{} \ill{} centre de $G_0$) \struck{\ill{}} \uncertain{Donc dans ce cas} \ill{}
$=$ \uncertain{ordre}], on voit
\marginal{en diagonale, à gauche, plusieurs lignes serrées où se lisent $G_0^{+}$,
$\mathbb{Z}/2\mathbb{Z}$, $D_2$, « Mais », To~3, « octaèdre », $\{1,\varepsilon\}$ ;
le reste \ill{}}

(49) $\left\{\begin{array}{l}
\text{a) ou bien } H \text{ \struck{\ill{}} est déterminé par son étoile,}\\
\text{b) ou bien } P_0 \text{ est un octaèdre (\uncertain{angulaire}), } Z'=\{1,\varepsilon\},
\ \varepsilon \text{ est}\\
\quad\text{l'antipodisme, i.e.\ l'unique élément}\\
\quad\text{central} \neq \mathrm{id} \text{ de } G_0=\mathrm{Aut}(P_0)\ 
(\simeq \mathfrak{S}_4\times\{\pm 1\}).
\end{array}\right.$

Et \uncertain{dans} \ill{} \ill{} que To~3, To~4, (\uncertain{pour que les sommets}
\ill{} \ill{} $=$ \uncertain{ordre}], alors $P_0$ est le \ill{} :
$Z'=\{1,\varepsilon\}$ si $n$ impair ou $n=4$, ]
$Z'\simeq\{1,\sigma\}\times\{1,\varepsilon\}$ dans le cas $n$ pair $\geqslant 6$.

\page{124}\note{p. 63 de l'auteur.}
Ainsi, la classe $\pi$ (se trouve être un espace
\add{de structures polyédrales « équivalentes »}
homogène sous le groupe \emph{commutatif} $Z'$),
il en résulte que le stabilisateur d'un $i\in\pi$
dans $Z'$ est indépendant de $i$, c'est donc aussi
le stabilisateur de $F_0$ \add{ou de $A_0$}, i.e.\ formé des
$\lambda\in Z'=\mathrm{Centr}_{\mathfrak{S}_S}(G_0)$ tels qu'on ait $\forall s,t\in S$
\[
(50)\qquad \{s,t\}\in A_0 \Longleftrightarrow \{s,\lambda t\}\in A_0
\]
ce qui implique d'ailleurs (remplaçant $\{s,t\}$ par $\{\lambda t,s\}$)
\[
\{\lambda t,\lambda s\}=\{s,\lambda t\}\ \ldots\in A_0 \Longleftrightarrow \{\lambda s,\lambda t\}\in A_0
\]
\note{la ligne de gauche porte une biffure illisible avant l'accolade.}
d'où
\[
\lambda\in G_0 \text{ et même } \lambda\in\mathrm{Cent}(G_0)\ (\text{\uncertain{même \ill{}}})
\]
puisque $\lambda$ centralise $G_0$. \struck{\ill{}} Mais \ill{} pour dire que (50)
signifie \ill{} :
\[
(51)\qquad \mathrm{Et}_{A_0}(t)=\mathrm{Et}_{A_0}(\lambda t)
\]
— ce qui implique que $\lambda=\mathrm{id}$, sauf dans le seul cas où \uncertain{les étoiles}
\uncertain{ne déterminent pas} les \uncertain{sommets}, i.e.\ le cas \uncertain{bien connu} de
l'octaèdre, \uncertain{si on utilise} \uncertain{tous} \ill{} que To~3, To~4 — et
celui de la \uncertain{double pyramide}. On trouve donc

\struck{$Z'$ \ill{} \ill{} \ill{} d'\uncertain{autant} de To~3, To~4, si
$\mathbb{Z}/\ill{}$ $=$ \ill{} \ill{} \ill{} conditions}

(52) \struck{\ill{}} $\pi$ est un torseur sous $Z'$, i.e.\ \uncertain{le} stabilisateur
\struck{\ill{}} $A_0$ dans $Z'$ est $\{1\}$, i.e.\ (50) implique
$\lambda=1$, sauf dans le cas \uncertain{double} où $P_0$ est une
double pyramide (49 b), où \ill{} $Z'\simeq G_0$, \struck{\ill{}}
\uncertain{où ce} stabilisateur est réduit à $\{1,\varepsilon\}$. \struck{\ill{}}
\marginal{en diagonale, à gauche : \uncertain{moyennant} To~3, To~4}

\page{125}
Dans le cas exceptionnel de la double
pyramide, \ill{} \ill{} $\pi$ réduit : un élément $A_0$
\struck{\ill{}} i.e.\ $Z'=\{1,\varepsilon\}$ \ill{} $n$ impair ou $n=4$,
et dans le cas \uncertain{inverse} ($n$ pair $\geqslant 6$) $\pi$ a exactement
deux éléments, qui correspondent
: la double pyramide sur les \ill{} \uncertain{ensembles}
$E$, \uncertain{munis} des deux structures \uncertain{polygonales}
\uncertain{affines} (d'ordre \struck{\uncertain{pair}} \add{$n\equiv 0\ (4)$ et} $\geqslant 8$) déduites l'une
de l'autre par « \uncertain{multiplication} par $\frac{n}{2}-1$ » …

J'\uncertain{laisse} \ill{} pour le reste en cas
de la double pyramide, \ill{} dans le cas
de l'octaèdre, i.e.\ pratiquement je \uncertain{donne}
que \struck{\uncertain{les \ill{}}} \add{\uncertain{seulement}} \ill{} tous même \uncertain{ordre} — dans
le seul cas exceptionnel où $\pi$ n'est
\uncertain{pas un} torseur sous \struck{$Z$} $Z'$ est \uncertain{dans} \ill{}
\ill{} $P_0$ est l'octaèdre, $\pi$ réduit : un
\ill{}, $Z'=\{1,\varepsilon\}$.
\note{ce paragraphe est marqué d'un long trait vertical dans la marge gauche.}

Notons que par transport de structure,
les permutations $f$ de $S$ qui transforment
un élément $F_0$ de $\pi$ en un autre $F_1\in\pi$,
\uncertain{transforment} \add{\uncertain{tout} \ill{}}
\ill{} $\pi$ en lui-même, et induisent
\uncertain{des} bijections de $\pi$ sur lui-même. Soit donc

\page{126}\note{p. 64 de l'auteur.}
\[
(53)\qquad G=\{f\in\mathfrak{S}_S \mid f \text{ prolongé à } \mathfrak{P}(\mathfrak{P}_2(S))\}
\]
\uncertain{applique} $\pi\subset\mathfrak{P}(\mathfrak{P}_2(S))$ en lui-même, donc satisfait
$f(\pi)=\pi$.
\marginal{en diagonale, à gauche : NB Si $\pi$ \uncertain{réduit} \ill{} $G=G_0$ \ill{}
\ill{} \ill{} ($\pi\geqslant 2$) \ill{} \ill{} \uncertain{triviale} \ill{} $Z'$.}

\struck{L'\ill{}} Notons qu'on a manifestement
\[
G_0 \overset{\mathrm{d\acute{e}f}}{=} \mathrm{Aut}(S,A_{i_0}) \subset G
\]
$\Vert$ \quad Stabilisateur de $i_0\in\pi$ dans $G$, opérant sur $\pi$

$=$ noyau de $G\to\mathfrak{S}_\pi$

(54) donc $G_0$ est un sous-groupe distingué de $G$, et on a
\[
(55)\qquad G/G_0 \hookrightarrow \mathfrak{S}_\pi .
\]
Mais $G$ normalisant $G_0$, normalise aussi son
centralisateur $Z'$, donc opère sur $Z'$ ;
et $G_0$ y opère trivialement, donc a priori
\[
(56)\qquad G/G_0 \longrightarrow \mathrm{Aut}_{\mathrm{gr}}(Z'),
\]
et il est clair par transport de structure
que (55) est compatible avec la structure
de $Z'$-espace homogène sur $\pi$, — en fait
le $Z'$-torseur $\pi$, \struck{\ill{}} i.e.\ on a a priori
\[
(57)\qquad G/G_0 \longrightarrow \mathrm{Aff}_{Z'}(\pi) .
\]

\page{127}
Ce serait le moment d'ailleurs de
noter que l'on a
\[
(58)\qquad Z' \subset G
\]
car on voit \struck{\ill{}} pour $(s,t)\in S^2$
\[
\{s,t\}\in A_{i_0} \Longleftrightarrow \{s,\lambda t\}\in A_{\lambda(i_0)}
\Longleftrightarrow \{\lambda s,\lambda t\}\in A_{\lambda.(\lambda i_0)},
\qquad \lambda.(\lambda i_0)=\lambda^2 i_0,
\]
donc
\[
(59)\qquad \underbrace{\lambda(A_{i_0})}_{\text{transport de structure}}
= A_{\lambda^2.(i_0)}
\]
où $\lambda^2.(i_0)$ désigne l'opération de $Z'$ sur $\pi$ définie par (45) p.~61,
et l'opération de $G$ sur $Z'$ est induite par
automorphismes intérieurs. Ceci dit, on \uncertain{va}
\uncertain{poser}
\marginal{en diagonale, à gauche : \uncertain{Attention}, \uncertain{\ill{}}
\uncertain{déterminant} \struck{\ill{}} \ill{} $\det(\lambda\,\mathrm{id}_M)=\lambda^2$
($\neq\lambda$ en \uncertain{général} !) (\ill{} par \ill{})}

(60) (To~7) $Z'$ \struck{\ill{}} est central dans $G$ (NB.\ il contient \struck{\ill{}}
le centre de $G$, donc To~7
signifie aussi
\[
Z'=\mathrm{Cent}(G)\ ),
\]
qui équivaut \uncertain{au} fait que $G$ opère
sur $\pi$ par translations, i.e.\ (56) trivial i.e.\ (57)
se \uncertain{factorise} par un homomorphisme de groupes
\[
(61)\qquad G \xrightarrow{\ \delta\ } Z',
\]
de sorte qu'on a finalement une suite
exacte canonique

\page{128}\note{p. 65 de l'auteur.}
(62)

\begin{tikzcd}
1 \arrow[r] & G_0 \arrow[r] & G \arrow[r, "\delta"] & Z' \\
 & & \mathrm{Cent}(G)=Z' \arrow[u] \arrow[ur, "\lambda\mapsto\lambda^2"'] &
\end{tikzcd}

\note{la flèche verticale relie $G$ et $\mathrm{Cent}(G)=Z'$ ; son sens est mal
marqué sur la page, on l'a orientée selon l'inclusion (58).}

On voit donc que $\delta(G)$ contient \struck{$Z'^{2}$}
\[
Z'^{2}=\{\lambda^2 \mid \lambda\in Z'\},
\]
mais cela ne signifie pas pour autant que
$\delta$ soit surjectif, i.e.\ que $G$ soit transitif
sur l'une des structures $\pi$. On voit de
suite que l'orbite d'un $i_0\in\pi$ sous $G$,
correspondant à la structure $A_{i_0}=A_0$ sur $S$,
est formée de l'ensemble des structures
$A_\lambda=A_{\lambda.i_0}$ sur $S$ « équivalentes » à $A_{i_0}$
(au sens de la p.~61) qui lui sont de
plus \emph{isomorphes}. L'axiome (To~5) — qu'on
n'avait pas utilisé jusqu'à présent \uncertain{ici})
signifie justement qu'il en \uncertain{trouve} ainsi
\uncertain{tout} \struck{\ill{}} $\pi$, i.e.\ que (61), est une
suite exacte \uncertain{courte}
\[
(63)\qquad 1\longrightarrow G_0\longrightarrow G \xrightarrow{\ \delta\ } Z' \longrightarrow 1 .
\]

\section{VIII Axiomatique : reformulation}

\page{129}\note{p. 66 de l'auteur ; le titre, souligné, est le sien.}
Une structure \emph{triangulaire} sur un ens $S$ est
\add{la donnée d'}un ens
\[
(1)\qquad F_0 \subset \mathfrak{P}_3(S)
\]
\add{un des faces} satisfaisant
\[
(2)\qquad S=\bigcup_{f\in F} f ,
\]
et on lui demande ici de satisfaire les
conditions

(3) (To~1) $F_0$ est l'ens des faces pour une \add{connexe}
structure de 2-polyèdre \struck{\ill{}} combinatoire
\uncertain{ayant} $S$ comme ens de sommets,
et comme ens des arêtes
\[
A_0=\{a\in\mathfrak{P}_2(S)\mid \exists f\in F_0,\ f\supset a\},
\]
et les relations d'incidence évidentes,

ce qui signifie que pour $\forall\, a\in\mathfrak{P}_2(S)$, s'il
existe $f\ni F_0$ \uncertain{avec} $f\supset a$, il en existe exactement
2, et que l'\uncertain{ensemble} les conditions de connexité
globale (deux \struck{\ill{}} \add{faces} de $S$ sont reliées par
une « chaîne de faces ») et locale en tout $s\in S$
(même chose, \uncertain{avec} des faces $\ni s$). On suppose
en plus que $F_0$ est déterminé par $A_0$ :
\[
(4)\qquad (\mathrm{To}\,2)\quad F_0=\{f\in\mathfrak{P}_3(S)\mid \forall a\in\mathfrak{P}_2(f),\ \text{on a } a\in A_0\}.
\]
Considérons maintenant les \struck{trois} \add{deux} groupes
\note{« $f\ni F_0$ » est écrit tel quel ; on attendrait $f\in F_0$.}

\page{130}
\[
(5)\qquad \text{\struck{\ill{}}}\ G_0=\mathrm{Aut}(S,F_0)=\mathrm{Aut}(S,A_0)=\mathrm{Aut}\,P_0
\]
(où $P_0=(S,A_0,F_0$, relation d'incidence …$)$)
\[
(6)\qquad \widehat{Z}=\mathrm{Centr}_{\mathfrak{S}_S}(G_0)
\]
\marginal{en diagonale, à gauche : je \uncertain{laisse} \uncertain{tomber} ici (5), (6)
\uncertain{finalement} !}
\struck{\ill{}} et soit
\[
(7)\qquad Z^{*}=\{\lambda\in\widehat{Z}\mid \exists f\in\mathfrak{S}_S \text{ t.q. }
\forall \{s,t\}\in S,\ \text{on a}\ 
\{s,t\}\in A_0 \Longleftrightarrow \{f^{-1}s,f^{-1}\lambda t\}\in A_0\}.
\]
\struck{i.e.}
On peut dire aussi que $\lambda\in\widehat{Z}$ est dans $Z^{*}$ ssi
il satisfait les deux conditions :

(8) $\left\{\begin{array}{l}
\text{a) } \forall (s,t)\in S,\ \{s,t\}\in A_0 \Longleftrightarrow \{\lambda^{-1}s,\lambda t\}\in A_0\\
\quad[\text{d'où un } A_\lambda\subset\mathfrak{P}_2(S), \text{ défini de façon unique}\\
\quad\ \text{par la condition que } \{s,t\}\in A_0 \Longleftrightarrow \{s,\lambda t\}\in A_\lambda]\\
\text{b) le graphe } \Gamma_\lambda=(S,A_\lambda) \text{ est isomorphe}\\
\quad\text{au graphe } \Gamma_0=(S,A_0)
\end{array}\right.$
\marginal{à droite de b), séparé d'un trait : cela implique que $\Gamma_\lambda$ est
\uncertain{associé} à une structure triangulaire $F_\lambda$ sur $S$.}

On suppose
\[
(9)\qquad (\mathrm{To}\,3)\quad \text{L'ensemble } Z^{*}\subset\widehat{Z}\subset\mathfrak{S}_S
\text{ est formé}
\]
\quad de permutations qui commutent
deux à deux, \struck{\ill{}}

On voit (en distinguant le cas où $H$
\uncertain{sommet} est déterminé par son étoile, et celui
où il n'en est rien i.e.\ où $P_0$ est une
double pyramide) que la condition (To~3)

\page{131}\note{p. 67 de l'auteur.}
\uncertain{équivaut} à : la \add{même} \uncertain{condition} \add{des} \ill{}
(10) $Z^{*}$ est un sous-groupe \struck{\ill{}} \add{$Z\subset\mathfrak{S}_S$},
\uncertain{sera} donc nécessairement commutatif.
\marginal{en diagonale, à gauche : NB Il faudrait \uncertain{simplifier} $Z^{*}$ \ill{}
(To~3') \ill{} \ill{} \ill{} des \uncertain{éléments} \ill{} \uncertain{double}
\uncertain{pyramide} \ill{} ($\mathrm{Aut}$ $\pi$ (p.~57) (35)), \uncertain{structure}
\ill{} \ill{} (11) \ill{}}

Introduisant maintenant
\[
(11)\qquad \pi \subset \mathfrak{P}(\mathfrak{P}_2(S)), \qquad
\pi=\{A_\lambda \mid \lambda\in Z\}
\]
\marginal{à droite, séparé d'un trait : Les groupes d'automorphismes des $\Gamma_\lambda=(S,A_\lambda)$
sont tous égaux à $G_0$, et $Z$ opère sur $\pi$ de façon transitive
($G_0$, \ill{} et \ill{} ne changent pas si on remplace $(S,F_0)$ par $(S,F_i)$ ($i\in\pi$)).}
\[
(12)\qquad G\subset\mathfrak{S}_S, \qquad
G=\{f\in\mathfrak{S}_S \mid f(A_0)\in\pi\}
\]
i.e.\ $f(\pi)=\pi$, \struck{on \ill{}} i.e.\ $\exists\lambda\in Z'$ t.q.\ $f(A_0)=A_\lambda$.
\marginal{en diagonale, à gauche : \struck{On aurait} \uncertain{présentation} \ill{}
par la suite $G$, on suppose $\pi$ non (12) réduit : un seul élément (\uncertain{sinon}
$G=G_0$, tout est trivial) (\ill{} \ill{}). Dans le cas de la double pyramide,
\struck{\ill{}} \add{en \uncertain{tous} cas} cela signifie que $n\geqslant 8$, $n\equiv 0\ (4)$
\struck{bipyramidale} et \ill{} dans le cas de l'\uncertain{orbite} des \ill{}
(13) \uncertain{stabilisateur} de $Z^{*}$ sur $\pi$ est $\{1,\varepsilon\}$.}

\struck{\ill{}} \struck{$A_i$, $F_i$} \add{\uncertain{notons}} pour tout $\lambda\in Z$,
\[
(13)\qquad A_i\in\mathfrak{P}(\mathfrak{P}_2(S)),\quad F_i\in\mathfrak{P}(\mathfrak{P}_3(S))
\]
les ens.\ d'arêtes, resp.\ de faces, associés à un
$i\in\pi$ \quad ($i=\lambda i_0$, $\lambda\in Z$). On \struck{suppose \ill{}} \add{note que} $G$
\note{sous $\lambda i_0$ : « il vaudrait mieux noter $T_\lambda i_0$ ».}
ne dépend que de $\pi\subset\mathfrak{P}(\mathfrak{P}_2(S))$ ou de
l'ens.\ correspondant $\pi'\subset\mathfrak{P}(\mathfrak{P}_3(S))$. On a
\[
(14)\qquad \mathrm{Cent}(G)\subset Z\subset G \supset G_0
\]
\struck{\ill{}} \struck{donc supposons}

avec $G_0$, $Z$ distingués dans $G$, et on a
\[
(15)\qquad G/G_0 \hookrightarrow \mathrm{Aff}_Z(\pi)
\]
$G/G_0$ opère \emph{transitivement} sur $\pi$.
\marginal{à droite de (15) : [\struck{\ill{}} Dans le cas bipyramidal avec $n\equiv 0\ (4)$,
$n\geqslant 8$, il faut remplacer $Z$ par $Z/\{1,\varepsilon\}$]}

On suppose

\page{132}
\[
(16)\qquad (\mathrm{To}\,4)\quad Z \text{ est central dans } G \text{ i.e.\ } Z=\mathrm{Cent}(G),
\]
i.e.\ l'opération de $G/G_0$ sur $\pi$ se fait
\struck{\ill{}} via un homomorphisme canonique
\[
G/G_0 \longrightarrow Z \quad [\text{resp.\ } Z/\{1,\varepsilon\} \text{ dans le cas bipyramidal}]
\]
\marginal{en diagonale, à gauche : NB Si l'on écrit (16) \ill{} \ill{} \ill{} \ill{},
\uncertain{nécessairement} $G'=\mathrm{Ker}(G\to\mathrm{Aut}(Z))=\mathrm{Centr}_G(Z)\supset Z$,
puis $Z'=\mathrm{Im}(G'\to Z)\supset Z^{2}$, puis la \uncertain{diagramme} canonique
$1\to G_0\to G'\to Z'\to 1$ \uncertain{surjectif} \ill{} $Z'\subset Z$ \ill{}
\ill{} $\pi'=\{A_\lambda\mid\lambda\in\pi\}$, \ill{} \ill{} \ill{} \ill{} \ill{} $Z$
\ill{} \ill{} (17) \ill{} \ill{} dans le cas bipyramidal, il faut remplacer $Z$ par
$Z/\{1,\varepsilon\}$ \struck{\ill{}}}

qui \uncertain{s'insère} dans une suite exacte
\[
(17)\qquad 1\longrightarrow G_0\longrightarrow G\xrightarrow{\ \chi'\ } Z\longrightarrow 1 ,
\]

\begin{tikzcd}[column sep=small]
 & G \arrow[r, "\chi'"] & Z \\
 & Z=\mathrm{Cent}(G) \arrow[u, hook] \arrow[ur, "\lambda\mapsto\lambda^2"'] &
\end{tikzcd}

qui est donc nécessairement \uncertain{canonique} :
\struck{\ill{}} $(S,F_0)$ satisfaisant (To~1) à (To~4), on
\uncertain{associe} : l'ens.\ $\pi$ de structures triangulaires sur
$S$ défini par \struck{$P_0$} $=(S,F_0)$.

Appelons les structures triangulaires satisfaisant
les axiomes (To~1) à (To~4) « \emph{prégroupales} \add{symplectiques} »
\add{\uncertain{triangulaires}}. On introduit sur l'ens des
\struck{\uncertain{structures}} structures triangulaires prégroupales sur
\struck{\ill{}} $S$ la relation d'équivalence
\marginal{en diagonale, à gauche : \uncertain{élégant} \ill{} \ill{} \ill{} la définition
\uncertain{des} (18), \ill{} \ill{} \ill{} \ill{} (To~4) \ill{} \uncertain{peut-être} \ill{} …}
\[
(18)\qquad \left\{\begin{array}{l}
(S,F_0) \text{ équivalente à } (S,F_1), \text{ \struck{si il existe} s'il existe}\\
\text{\struck{$\lambda$}}\ \lambda\in Z_{F_0} \text{ tel que } F_1=\lambda(F_0), \text{ i.e.\ \struck{s'il existe}}\\
\quad \text{a) } (S,F_1) \text{ isom.\ à } (S,F_0)\\
\quad \text{b) } \exists\lambda\in\mathfrak{S}_S, \text{ telle que } \lambda \text{ commute à}\\
\qquad G_0=\mathrm{Aut}(S,F_0), \text{ et que l'on ait}\\
\qquad \{s,t\}\in A_{F_0} \Longleftrightarrow \{s,\lambda t\}\in A_{F_1}.
\end{array}\right.
\]
\struck{Soit} Une structure triangulaire prégroupale

\page{133}\note{p. 68 de l'auteur.}
sur \add{droite …} un $S$ sera par définition, la donnée
d'une classe d'équivalence$_\pi$, au sens de (18)
de structures tr.\ sympl.\ prégroupales.
\marginal{en diagonale, à gauche : $\pi\subset\mathfrak{P}(\mathfrak{P}_3(S))$
ou $\pi\subset\mathfrak{P}(\mathfrak{P}_2(S))$}
À une telle structure est donc associé
une suite exacte canonique (17), où
$G_0$ est le groupe d'automorphismes commun des
structures $(S,F_i)_{i\in\pi}$, et où $Z$ est le
groupe $Z_{F_i}$, commun pour tous les $i\in\pi$,
$\pi$ étant un torseur canonique sous
$Z$.

\emph{NB} À vrai dire, cette axiomatique n'est pas
encore réduite au minimum ; je pense, en un
sens, que les structures triangulaires \uncertain{ici}
\uncertain{pourront} qu'en vertu logique « \uncertain{surjectives} » —
les réflexions aboutissant aux axiomes
To~3, To~4 (To~1 et To~2 \uncertain{étant} \uncertain{naturels}) pourraient se dissiper
par le plus naturellement dans le contexte des
\emph{graphes} plutôt que dans celui des structures
triangulaires qui leur sont associées : \emph{certains} graphes.

Pour \uncertain{ces} raisons, qui ont \uncertain{à} aux structures
triangulaires, on voit que la structure
triangulaire prégroupale
\[
(19)\qquad \pi\in\mathfrak{P}(\mathfrak{P}(\mathfrak{P}_3(S)))
\]

\page{134}
peut aussi se décrire entièrement : \uncertain{l'aide}
de
\[
(20)\qquad F=F_\pi=\bigcup_{i\in\pi}F_i=\{f\in\mathfrak{P}_3(S)\mid \exists i\in\pi \text{ tel que}
\]
\quad $f$ soit une face pour la structure polyédrale $P_i\}$,

— en proposant d'expliciter les conditions
sur $F$, pour que $(S,F)$ provienne d'une
« structure triangulaire prégroupale », par
\uncertain{sous} les mêmes \emph{uniques}, les $F_i$ s'interprétant
alors comme les ens.\ des sommets des
composantes connexes du graphe
\[
(21)\qquad \Gamma_F=\{F,A,R\}
\]
\marginal{en diagonale, à gauche : $R$ la relation d'incidence}
où on définit
\[
(22)\qquad A=\{a\in\mathfrak{P}_2(S)\mid \exists f\in F,\ f\supset a\}
\]
$A\subset\mathfrak{P}_2(S)$, (contenu de $F$, comme
$A=\bigcup_{i\in\pi}A_i$) :
\note{la phrase « $A\subset\mathfrak{P}_2(S)$, contenu de $F$ … » est écrite en surcharge
à droite de (22) ; l'ordre en est incertain.}

(en fait, \ill{} \ill{}
\[
(23)\qquad \pi\simeq\pi_0(\Gamma_F).
\]
Les axiomes \struck{sont \quad $\bigcup_{f\in F}f=S$ \quad (pour \ill{} \ill{} \ill{}
structure triangulaire)}
\[
(\mathrm{T}\,1)\qquad \forall a\in A, \text{ existent exactement \emph{deux} } f\in F
\]
\quad qui contiennent $a$,

ce qui permet de considérer $\Gamma_F$ comme
un graphe. Prenant les composantes connexes

\page{135}\note{p. 69 de l'auteur.}
de ce dernier, on trouve une décomposition
\[
(24)\qquad F=\coprod_{i\in\pi}F_i, \qquad A=\coprod_{i\in\pi}A_i
\]
les $F_i$, $A_i$ \uncertain{non vides} et mutuellement disjoints.
\note{le numéro (24) est écrit sur un autre, biffé.}
Les \uncertain{exigeons} alors (cf.\ p.~53)

(25) (T~3) $\forall i\in\pi$, $(S,F_i)$ est une structure
\emph{triangulaire stricte} $P_i$, i.e.\ $\bigcup_{f\in F_i}f=S$, et
on a connexité \uncertain{du contour local en}
$s$ défini par $P_i$,

cf.\ T3 a) b) p.~53' pour une explicitation.
\note{« p. 53 », « p. 53' » : pagination de l'auteur, en amont de ce lot.}

Les autres axiomes pouvant s'exprimer, en
disant que la famille $(F_i)$ de structures
triangulaires sur $S$ \struck{\ill{}} \uncertain{provient} exactement d'une \uncertain{donnée}
\uncertain{d'équivalence} \ill{} de structures triangulaires
\uncertain{sympl.}\ symplectiques prégroupales.

Je \uncertain{faire} en \uncertain{arrivant} enfin, progressivement
des conditions plus strictes encore,
\uncertain{qui} un \uncertain{arrivera} enfin \uncertain{aux} structures « \uncertain{régulières} »
\struck{\ill{} \ill{} \ill{} \ill{}} « \uncertain{régulières} » \ill{}
\uncertain{triangulaires}, celles déduites des \uncertain{modules}
\uncertain{libres} \uncertain{monogènes} de rang 2 sur des $\mathbb{Z}/n\mathbb{Z}$. Pour
l'instant, l'attention principale \struck{\ill{}} \uncertain{sera}
sur les phénomènes \ill{} au cas d'un $\pi$ non
réduit : un point, i.e.\ d'un $Z$ qui n'est
pas réduit : $\{1\}$ (ou à $\{1,\varepsilon\}$, dans le cas
\marginal{en diagonale, à gauche, en plusieurs lignes serrées : En \uncertain{dimension} 2,
\ill{} $(S,F)$ \ill{} (25 bis) \ill{} $G$ \ill{} \ill{} $(S,F)$ \ill{} \ill{} \ill{} \ill{}
$\overline{A}$, $A$, $S$ \ill{} \ill{} $\mathbb{Z}/3\mathbb{Z}$ (\uncertain{canonique}),
$\mathbb{Z}/12\mathbb{Z}$, $D_{13}$ \ill{} \ill{} \ill{} $\mathbb{Z}/13\mathbb{Z}$ \ill{} ;
lecture très incertaine.}

\page{136}
bipyramidal). Je \uncertain{voudrais} \uncertain{toujours} \add{\ill{}}
\uncertain{définir} un isomorphisme
\[
(26)\qquad Z \xrightarrow{\ \sim\ } (\mathbb{Z}/n\mathbb{Z})^{*}/\{\pm 1\}
\]
\quad (resp.\ $Z/\{1,\varepsilon\}$)

pour $n$ convenable.

On va donc revenir, pour fixer les
idées, à la donnée d'une structure triangulaire
symplectique prégroupale. On supposera
\[
(27)\qquad (\mathrm{To}\,5)\quad \text{Les sommets de } S \text{ sont tous de \struck{\ill{}} ordre fini,}
\]
\quad soit $n\in\mathbb{N}$ ($n\geqslant 3$).

Soient maintenant $\lambda\in Z$, $s\in S$, on a
par définition de $A_\lambda$
\[
(28)\qquad \mathrm{Et}_{A_0}(s)=\mathrm{Et}_{A_\lambda}(\lambda s)
\]
\note{sous la formule : « $E_0$ » sous le premier membre, « $A_0$ » sous le second, ajouts de sa main.}
donc l'un des étoiles \add{pour $A_0$} \struck{des points} \add{des} sommets $s\in S$
est la même que l'une \uncertain{des} $E_\lambda$ pour $A_\lambda$ — simplement
les « contours » de ces étoiles changent, quand on
change de structure $A_i$ ($i\in\pi$) — et bien entendu,
aussi, la structure polygonale.

Notons qu'une structure triangulaire $(S,F_0)$
satisfaisant \struck{\ill{}} (To~2) est déterminée par la
connaissance de l'ens $\mathcal{E}$ des étoiles qu'elle
définit, ici un ensemble de (27)
\[
(29)\qquad \mathcal{E}_0\subset\mathfrak{P}_n(S),
\]
\struck{$E\in\mathcal{E}_0 \Rightarrow \mathrm{Card}\,E=n$} (\ill{}
et \uncertain{avec} la donnée \uncertain{plus} \uncertain{précise} \struck{\ill{} \ill{}}

\page{137}\note{p. 70 de l'auteur. La numérotation passe de (29), p.~136, à (33) : les formules (30) à (32) sont au feuillet suivant, p.~138, qui fait suite à la p.~136.}
En effet, \struck{pour \ill{} \ill{} structure \ill{} polygonale}
\add{\uncertain{voyons} d'abord les structures}
\struck{\ill{} \ill{}}
\ill{} polygonales orientées, sont $\overrightarrow{\mathrm{Pol}}(E)$,
alors \ill{}.

(33) Soit $E$ de cardinal $n$, dans $(\mathbb{Z}/n\mathbb{Z})^{*}$ opère
canoniquement sur $\overrightarrow{\mathrm{Pol}}(E)$.

En effet, une structure de $\overrightarrow{\mathrm{Pol}}(E)$ s'interprète
\uncertain{précisément} comme une opération de $\mathbb{Z}/n\mathbb{Z}=\Lambda_n$ sur $E$
faisant de $E$ un $\Lambda_n$-torseur
(lire une « droite affine sur $\Lambda_n$ »)
donc $\mathrm{Aut}(\Lambda_n)$ opère sur l'ens de ces
structures, par composition. Si on identifie
un $\omega\in\overrightarrow{\mathrm{Pol}}(E)$ comme une permutation
circulaire, on trouve
\[
(34)\qquad \underbrace{T_\nu\,\omega}_{\text{effet sur } \omega\in\overrightarrow{\mathrm{Pol}}(E)\subset\mathfrak{S}_E \text{ du } T_\nu \text{ associé à } \nu\in(\mathbb{Z}/n\mathbb{Z})^{*}}=\omega^{\nu}
\]
On a donc
\[
(35)\qquad \left\{\begin{array}{l}
T_{-\nu}\,\omega=\omega^{-\nu}=(\omega^{\nu})^{-1}=(T_\nu\,\omega)^{-1}\\
\qquad \text{structure polygonale orientée opposée : } T_\nu\,\omega\\
T_\nu\,\omega^{-1}=(T_\nu\,\omega)^{-1}
\end{array}\right.
\]
i.e.\ $T_\nu$ commute à \struck{\ill{}} $\omega\mapsto\omega^{-1}$
passage d'une str.\ de polyg.\ orient.\ à l'opposée.

\page{138}\note{suite de la p.~136 (formules (30) à (32)), insérée avant la p.~137 dans la lecture.}
\[
(30)\qquad \forall E\in\mathcal{E}_0,\ \text{\struck{\ill{}}}\ \text{la donnée de la structure de}
\]
\quad $n$-polygone combinatoire, ayant $E$ comme
ens.\ de ses sommets, i.e.\ $A_E\subset\mathfrak{P}_2(E)$

puis, qu'on récupère $A_0$ à partir de (29) (30)
par
\[
(31)\qquad A_0=\{a\in\mathfrak{P}_2(S)\mid \exists E\in\mathcal{E}_0 \text{ t.q.\ } a\in A_E\}.
\]
À vrai dire, pour \uncertain{dans} le cas où l'axiome (To~1) est
satisfait, il existe exactement \emph{deux} $s\in S$
tels que $\mathrm{Et}_{A_0}(s)$ contienne $a$ (donc exactement
\uncertain{deux} comme ceci) et aussi dans le cas octaédral,
ces étoiles associées sont distinctes, donc aussi dans
le cas de l'octaèdre, il y a aussi dans (31),
exactement deux $E\in\mathcal{E}_0$ telles que $E\supset a$, i.e.\
$a\in A_E$. (Dans le cas de l'octaèdre, par contre,
le $E$ en question est unique.)

Pour déterminer les structures $A_\lambda$ ($\lambda\in Z$),
connaissant \add{l'une $\mathcal{E}_\lambda$ des} ses étoiles \uncertain{est} égal : $\mathcal{E}_0$, il suffit
de savoir décrire, \struck{\ill{}} pour tout $E\in\mathcal{E}_0$, la
structure $n$-polygonale $A_{E,\lambda}\in\mathfrak{P}_2(E)$ induite sur
\struck{\ill{}} $E$ par $\Gamma_\lambda=(S,A_\lambda)$. Notons maintenant
le fait mirifique
\[
(32)\qquad \left|\ \begin{array}{l}
\text{Soit } E \text{ un ens.\ fini, de cardinal } n\geqslant 3. \text{ Alors}\\
\text{le groupe } Z'_n=(\mathbb{Z}/n\mathbb{Z})^{*}/\{\pm 1\} \text{ opère canoniquement}\\
\text{sur l'ens.\ } \mathrm{Pol}(E) \text{ des structures polygonales}\\
\text{sur } E \text{ comme ens.\ de sommets.}
\end{array}\right.
\]
\marginal{en diagonale, à gauche de (32) : on voit de suite que cette opération est fidèle.}

\page{139}\note{p. 71 de l'auteur. La seconde formule (33) de la page reprend un numéro déjà employé p.~137.}
donc l'opération de $(\mathbb{Z}/n\mathbb{Z})^{*}$ sur $\overrightarrow{\mathrm{Pol}}(E)$
passe au quotient en une opération sur
\[
\text{\struck{$\mathrm{Pol}(E)$}}\qquad \mathrm{Pol}(E)=\overrightarrow{\mathrm{Pol}}(E)/(\omega\mapsto\omega^{-1}),
\]
\uncertain{enfin} l'élément $-1\in(\mathbb{Z}/n\mathbb{Z})^{*}$ opère
trivialement sur $\mathrm{Pol}(E)$, en vertu de (35),
d'où l'opération (32).
\marginal{en diagonale, à gauche : si $A\subset\mathfrak{P}_2(E)$ définit
la structure polygonale $\alpha\in\mathfrak{P}_2(E)$, on désigne par
$A^{\nu}$ la transformée par $\nu\in Z'_n$ \ill{}}

Pour $\nu\in(\mathbb{Z}/n\mathbb{Z})^{*}/\pm 1=Z'_n$, et
\uncertain{$\alpha$} structure polygonale sur un ens $E$, on
\uncertain{désigne} par $\alpha^{\nu}$ la structure \struck{\ill{}} \add{transformée}
par $\nu$. \uncertain{Ceci} fixé, j'énonce dès
l'axiome suivant :
\[
(33)\qquad (\mathrm{To}\,6)\quad \forall\lambda\in Z,\ \text{\struck{\ill{}}}\ \exists\,\nu\in\underbrace{(\mathbb{Z}/n\mathbb{Z})^{*}/\pm 1}_{Z'_n}
\]
\quad (évidemment unique) tel que
\quad pour tout $s\in S$, la structure
\quad polygonale \struck{induite} $\alpha_{\lambda.i_0,\lambda.s}$ induite
\quad par $A_{\lambda.i_0}=A_\lambda$ sur $\mathrm{Et}_{A_0}(s)=\mathrm{Et}_{A_\lambda}(\lambda s)$,
\quad soit égale à $(\alpha_{i_0,s})^{\nu}$.
\note{sous « $\lambda.i_0$ » : « plutôt : $T_\lambda i_0$ ».}

Cela détermine donc la structure $A_\lambda$ en
termes de $\nu=\nu(\lambda)\in Z'_n$, \add{notons-la $A^{(\nu)}$,} donc dans le cas
où $P_0$ n'est pas l'octaèdre (\add{\uncertain{puisque}} \uncertain{exception} \uncertain{des} \uncertain{triviaux} \uncertain{fixes}
$Z=\{1,\varepsilon\}$ opère trivialement sur $\pi$, réduit : un pt !),
cela détermine aussi $\lambda$ en termes de $\nu$ :

\page{140}
\[
(34)\qquad \left|\ \begin{array}{l}
\text{\struck{\ill{}}}\ \lambda s \text{ est l'unique élément de } S, \text{ dont}\\
\text{l'étoile pour } A_\lambda=A_0^{\nu} \text{ est égale à } \mathrm{Et}_0(s).
\end{array}\right.
\]
[\emph{NB} Dans le cas de l'octaèdre, rappelons que $Z=\{1,\varepsilon\}$,
et que pour $s\in S$, $\varepsilon(s)$ pour ces conditions
\uncertain{correspond} l'unique $s'\neq s$ qui ait même
étoile que $s$.] Ainsi on trouve un
\emph{homomorphisme injectif}
\[
(35)\qquad Z\hookrightarrow\Lambda'_n
\]
\quad (cas différent du cas octaédral, qui \uncertain{a} $\Lambda'_n=\{1\}$, $Z=\{1,\varepsilon\}$)

dont l'image est formée d'éléments $\nu$ satisfaisant
\uncertain{à} la condition suivante :

(36) L'ensemble \struck{$A_0^{(\nu)}=\{a\in\mathfrak{P}_2(S)\mid \exists E\ldots\}$}
\[
A_0^{(\nu)}=\coprod_{E\in\mathcal{E}_0}A_E^{(\nu)}
\]
\quad est l'ensemble des arêtes pour une
\quad structure de 2-polyèdre sur $S$ [ayant
\quad mêmes étoiles que $P_0$, et en fait, plus précisément,
\quad \emph{équivalente} à celle de $P_0$ au sens de (18)].
\note{les numéros (34), (35) de cette page reprennent ceux de la p.~137 ; les numéros sont ceux de l'auteur.}

Mais le \uncertain{simple fait ici}, \struck{\ill{}} pour une
structure de 2-polyèdre triangulée $P_0=(S,A_0,F_0)$
satisfaisant [To~1 à] To~2, \add{(To~5)}, que l'ens $A_0^{(\nu)}\subset\mathfrak{P}_2(S)$
défini sur $S$ une structure de graphe
provenant, elle aussi, d'une telle structure
triangulée — et ceci pour \emph{tout} $\nu\in\Lambda'_n$ —
\struck{\ill{}} est une propriété remarquable
et pourrait être très forte, qui mérite d'être examinée
\marginal{à gauche, en diagonale : cf.\ To~7 plus bas !}
\note{la phrase se poursuit au-delà de ce lot ; le symbole avant $\Lambda'_n$ est raturé et illisible.}

\end{document}
